\subsection{The exact finite-character obstruction to one initialization}

\begin{statement}{A.85}{Proved}
{A necessary and sufficient finite-character test}
Let $4\mid n$ and let the old cover $C$ and new puncture $c$
be as in (A.83.1)--(A.83.2), without the squareclass restriction
of A.83.  Let $B_c$ be the common abelian fibre field computed
by A.57 for
\[
 \mathcal E_c=\Q[Z_1^{\pm1},\ldots,Z_k^{\pm1}]
                         /(Z_i^n-(c-c_i)/a_i).
 \tag{A.85.1}
\]
Choose a finite set $S$ containing the ramified primes of its
splitting field, the coefficient, denominator and difference
primes specified in A.83, and all primes up to the uniform
augmented-cover bound (A.83.5).
For every $v\in S\cup\{\infty\}$ prescribe a point
$P_v\in C(\Q_v)$ with $T(P_v)\ne c$.
Define the finite character obstruction by the actual local
reciprocity maps:
\[
 \mathfrak o_c((P_v))=
   \prod_{v\in S\cup\{\infty\}}
      \operatorname{rec}_v(T(P_v)-c)\bigm|_{B_c}
                  \quad\in\operatorname{Gal}(B_c/\Q).
 \tag{A.85.2}
\]
The following two statements are equivalent:
\begin{enumerate}
\item $\mathfrak o_c((P_v))=1$.
\item There is $a\in\Q^\times$ such that the augmented cover
      $T-c=aY^n$ over $C$, with $Y\ne0$, has points at all
      places of $\Q$, and every prescribed $P_v$ lifts through
      its actual forgetful map.
\end{enumerate}
When these statements hold, $a$ can be chosen as a rational
$S$-unit times a product of distinct fresh primes whose
Frobenius fixes an old-fibre root tuple at $c$.  It can have
valuation one at a fresh prime outside any prescribed finite set.

Equivalently, the complete test consists of the finitely many
equalities
\[
 \sum_{v\in S\cup\{\infty\}}
   \operatorname{inv}_v(T(P_v)-c,\chi)=0
       \quad
 \bigl(\chi\in
       \operatorname{Hom}(\operatorname{Gal}(B_c/\Q),\Q/\Z)\bigr).
 \tag{A.85.3}
\]
Generators of this finite character group suffice.
All local points outside $S$ have zero evaluation in this
test.  Thus the obstruction depends on finitely many local
neighborhoods and cannot be removed by changing only the
local points outside $S$.
\end{statement}
\begin{proof}
Write $n=2^r v$ with $v$ odd and $r\ge2$.
By the actual field inclusion (A.57.5),
$B_c$ lies in a compositum of a multiquadratic field and
$\Q(\mu_{2^{r+1}})$.
The Galois group of the latter field is
$C_2\times C_{2^{r-1}}$.
Consequently every character of $B_c$ has order dividing
$2^{r-1}$, and in particular
\[
                              \operatorname{ord}(\chi)\mid n.
 \tag{A.85.4}
\]
The local invariant map sends the cyclic algebra
$(u,\chi)$ to $\chi(\operatorname{rec}_v u)$, in the
reciprocity convention specified in A.83.
Characters separate the finite abelian group
$\operatorname{Gal}(B_c/\Q)$, proving the equivalence of
(A.85.2) and (A.85.3).

\emph{Every outside-$S$ evaluation is zero.}
Fix a prime $p\notin S$ and an old local point
$P=(T,X_1,\ldots,X_k)$ with $T\ne c$.
All coefficients $a_i$ and differences $c-c_i$ are units,
all punctures are integral, $p\nmid n$, and $B_c/\Q$
is unramified at $p$.
For an unramified character the local cyclic invariant is
\[
              v_p(T-c)\,\chi(\operatorname{Frob}_p).
 \tag{A.85.5}
\]
We check every possible valuation case.

If $v_p(T)<0$, then
$v_p(T-c_i)=v_p(T)$ for every $i$.
The old equations give
$v_p(T)=n\,v_p(X_i)$, because $a_i$ is a unit.
Also $v_p(T-c)=v_p(T)$.
Equation (A.85.4) therefore makes (A.85.5) zero.

If $T$ is integral and $T-c$ is a unit, its valuation
is zero and the conclusion is immediate.
The remaining case is $T$ integral with $v_p(T-c)>0$.
Then every $X_i$ is a unit and reduction gives
\[
                 \bar X_i^{\,n}=(c-c_i)/a_i
                                         \quad\text{in }\mathbf F_p.
 \tag{A.85.6}
\]
This is a point of the actual finite etale fibre algebra
at $c$ over $\mathbf F_p$.
The Frobenius permutation of its geometric roots therefore
fixes a root tuple.
The character $\chi$ vanishes on the stabilizer of every
such tuple, by the definition and the proved character
calculation of $B_c$ in A.57.
Thus $\chi(\operatorname{Frob}_p)=0$, again giving zero in
(A.85.5).
This also covers all integral points reducing to another
old puncture, since then $T-c$ is a unit.
We have proved the claimed outside-$S$ assertion for every
old point where the cyclic expression is evaluated.

\emph{Necessity for any global coefficient.}
Suppose a rational $a$ as in the second assertion exists.
At $v\in S\cup\{\infty\}$, a lift of the prescribed point
gives
$T(P_v)-c=aY_v^n$.
At every other place choose any local point of the
augmented cover; it satisfies the same identity with a
nonzero $Y_v$.
By (A.85.4), the local cyclic classes of $a$ and
$T(P_v)-c$ are equal.
The latter evaluations are zero outside $S$, by the case
analysis just proved.
Global reciprocity for the constant cyclic algebra
$(a,\chi)$ now gives
\[
 0=\sum_v\operatorname{inv}_v(a,\chi)
   =\sum_{v\in S\cup\{\infty\}}
             \operatorname{inv}_v(T(P_v)-c,\chi).
 \tag{A.85.7}
\]
This proves necessity for every rational coefficient,
with no restriction on its prime factors.

\emph{The ray class determined by the obstruction.}
For sufficiency put
$\alpha_v=T(P_v)-c$ and use the same explicit $S$-unit as
in A.83:
\[
 d=\operatorname{sgn}(\alpha_\infty)
                      \prod_{p\in S}p^{v_p(\alpha_p)},\qquad
 u_p=\alpha_p/d\in\Z_p^\times .
 \tag{A.85.8}
\]
Choose $N_p>2v_p(n)$ so large that, for
$M=\prod_{p\in S}p^{N_p}$ and $R=\Q(\mu_M)$, one has
$B_c\subset R$.
This is an explicit finite choice: A.57 places $B_c$
inside $\Q(\mu_{2n},\sqrt P:P\mid\prod_i D_i)$, so the
cyclotomic modulus and the finitely many quadratic
conductors suffice.
Let $\rho\in(\Z/M)^\times$ be the unique CRT class with
$\rho\equiv u_p\pmod {p^{N_p}}$ for every $p\in S$.

Outside $S$, the number $d$ is a unit and every character
of $B_c$ is unramified.  Global reciprocity for $d$ and
the equality $\alpha_\infty/d>0$ give the exact relation
\[
 \prod_{p\in S}\operatorname{rec}_p(u_p)\bigm|_{B_c}
       =\mathfrak o_c((P_v)).
 \tag{A.85.9}
\]
Local cyclotomic reciprocity identifies the product on
the left, before restriction to $B_c$, with the automorphism
$\rho^{-1}$ of $R$: a local unit acts by its inverse on
the $p$-power roots and trivially on the prime-to-$p$ roots.
Thus
\[
              \rho^{-1}\bigm|_{B_c}=\mathfrak o_c((P_v)).
 \tag{A.85.10}
\]
Under the assumed vanishing, $\rho$ belongs to
$\operatorname{Gal}(R/B_c)$.
This is the exact ray-class condition; no larger splitting
condition has been imposed.

\emph{Generated stabilizers realize exactly that ray class.}
Let $L$ be the splitting field of (A.85.1), let
$G=\operatorname{Gal}(LR/\Q)$, and let $\Omega$ be its
root-tuple set.
For the restriction map
$\pi_R:G\to\operatorname{Gal}(R/\Q)$ and the actual
stabilizers $G_\omega$, Galois correspondence gives
\[
 R^{\langle\pi_R(G_\omega):\omega\in\Omega\rangle}
       =R\cap\bigcap_{\omega\in\Omega}\Q(\omega)=B_c.
 \tag{A.85.11}
\]
The last equality holds because this intersection inside
the abelian field $R$ is the largest common abelian
subfield of the root factors, and $B_c\subset R$.
It follows that the stabilizer images generate precisely
$\operatorname{Gal}(R/B_c)$.

Write $\rho=\gamma_1\cdots\gamma_t$ with each
$\gamma_j=\pi_R(g_j)$ for some $g_j$ fixing a root tuple.
Each image is a subgroup, so inverses are also allowed
factors.  Chebotarev supplies distinct arbitrarily large
primes $q_j\notin S$ whose Frobenius classes in $G$
are those of $g_j$.
The restriction to $R$ is unchanged by conjugation, and
the property of fixing some tuple is conjugacy invariant.
Hence the product of these primes has ray class $\rho$,
and each prime has an old residue-fibre point at $c$.
Choose additionally a distinct prime $q_*$ splitting
completely in $LR$, avoiding every prescribed finite set,
and put
\[
                          a=d\,q_*\prod_{j=1}^t q_j.
 \tag{A.85.12}
\]
The primes can all exceed the uniform bound in the
definition of $S$.
Here $q_*\equiv1\pmod M$ and $v_{q_*}(a)=1$.

The remaining local construction is explicit.
At $p\in S$ one has
$a/\alpha_p\in1+p^{N_p}\Z_p$, an $n$th-power subgroup
by Hensel with $N_p>2v_p(n)$.
Thus the given old coordinates of $P_p$ extend by an
$n$th root $Y_p$ of $\alpha_p/a$.
At the real place this ratio is positive.
At any fresh coefficient prime $q$, take $Y=1$ and
$T=c+a$; the residue old-fibre point lifts independently
in each $X_i$, since $q\nmid n$ and the derivatives are
units, exactly as in (A.83.16).
At every other prime the coefficients and puncture
differences are units, and the uniform genus and boundary
bound (A.83.4)--(A.83.5) supplies an interior finite-field
point, followed by smooth Hensel lifting.
This proves sufficiency and the assertion about the
prime factors of $a$.

Finally each character in (A.85.3) is continuous with
finite image.  Its value on $T(P_v)-c$ is locally
constant on points avoiding $c$.
The obstruction therefore is constant on a product
of sufficiently small neighborhoods of the finitely
many prescribed points.
The outside-$S$ calculation proves that all other local
choices have zero contribution, as asserted.
\end{proof}

\paragraph{The two earlier obstructions persist in finite neighborhoods.}
In A.69 and A.73, use their stated finite exponent $n$,
old global coefficients, rational base point $R$, and new
puncture $c$.
Choose $S$ as above, enlarging it to contain the ramification
of the constant cyclic classes used below, and prescribe the exceptional point
at two and the base point at every other place in
$S\cup\{\infty\}$.
The normalized cyclic expression used there is
$((T-c)/(R-c),\chi)$.
Its quotient by $(T-c,\chi)$ is one global constant,
whose total invariant sum is zero; all its finite
ramification is already included in $S$.
The explicit evaluations in A.69 and A.73 consequently
give
\[
                         \chi(\mathfrak o_c((P_v)))=\tfrac12.
 \tag{A.85.13}
\]
By necessity, there is no augmented all-place cover
with those old local points, regardless of the new
coefficient and regardless of the choices at all
primes outside $S$.
The same conclusion holds on sufficiently small
neighborhoods of the prescribed points.

If an ordinary completed old tuple has the given old
classes modulo $n$ and has its old local lifts in these
neighborhoods, any global completed class for the new
coordinate would project modulo $n$ to a rational
coefficient satisfying the forbidden second assertion.
That projection is the exact obstruction map.
A.90 constructs such a completed old tuple for the
fixed A.69 neighborhoods.
A.76 excludes the particular tuples fixed equal to
the rational base point on a density-one set; it does
not exclude every completed tuple in these finitely
many neighborhoods.
The obstruction there therefore persists for actual
completed old coordinates.

All group and prime-product calculations in this proof
are the exact maps certified in
\path{code/fibre_squareclass_certificate.py} and
\path{code/puncture_initialization_certificate.py}.
The general character group is computed by
\path{code/fibre_character_certificate.py};
the nonzero evaluations in (A.85.13) are the retained
certificates of A.69 and A.73.
